% Part: incompleteness
% Chapter: introduction
% Section: definitions

\documentclass[../../../include/open-logic-section]{subfiles}

\begin{document}

\olfileid{inc}{int}{def}

\olsection{تعريفات}

أول ما يقتضيه مشروع هيلبرت، في صوغ الرياضيات صوريًا وإثبات
اتساقها وتمامها، تعيين لغة وإطار منطقي وبديهيات. فلنفرض هنا
أن الرياضيات تقبل الصوغ بلغة من الرتبة الأولى؛ أي إن
ال!!{constant}s وال!!{function}s وال!!{predicate}s المناسبة،
مع الروابط والكمّات، تكفي للتعبير عن قضاياها. ومن المتفق عليه
على نطاق واسع أن لغة نظرية المجموعات تفي بذلك، وليس فيها
من الرموز غير المنطقية إلا $\in$. وقد تستبعد لأول وهلة
أن تتسع لغة بهذا الإيجاز لكل ذلك؛ وإثبات أنها تستوعب
الرياضيات كلها تقريبًا ثمرة عمل كثير. ولنقصر النظر الآن،
تيسيرًا، على الحساب: علم الأعداد الطبيعية~$\Nat$.
واللغة الملائمة له هي~$\Lang L_A$. وفي $\Lang L_A$
!!{predicate} ثنائي واحد $<$، و!!{constant} واحد
$\Obj 0$، و!!{function} أحادي واحد $\prime$،
و!!{function}s ثنائيتان $+$ و~$\times$.

\begin{defn}
مجموعة ال!!{sentence}s~$\Gamma$ تسمى \emph{نظرية} متى
كانت مغلقة تحت الاستلزام؛ أي متى تحققت المساواة
$\Gamma = \Setabs{!A}{\Gamma \Entails !A}$.
\end{defn}

ولتعيين النظريات طريقان يسيران. أولهما أن نجمع
ال!!{sentence}s الصادقة في !!{structure} معينة.
فمن أمثلته ال!!{structure} للغة $\Lang L_A$ التي
!!{domain}ها~$\Nat$، وتؤوّل فيها الرموز غير المنطقية
بتفسيراتها المعتادة.

\begin{defn}\ollabel{def:standard-model}
نسمي \emph{النموذج القياسي للحساب} ال!!{structure}~$\Struct N$
التي نعيّنها هكذا:
\begin{enumerate}
\item $\Domain N = \Nat$
\item $\Assign{\Obj 0}{N} = 0$
\item $\Assign{\Obj \prime}{N}(n) = n + 1$ لكل $n \in \Nat$
\item $\Assign{\Obj +}{N}(n, m) = n + m$ لكل $n, m \in \Nat$
\item $\Assign{\Obj \times}{N}(n, m) = n\cdot m$ لكل $n, m \in \Nat$
\item $\Assign{\Obj <}{N} = \Setabs{\tuple{n, m}}{n \in \Nat, m \in
  \Nat, n < m}$
\end{enumerate}
\end{defn}

وينبغي الفصل بين $\times$ و$\cdot$. فأما $\times$
فمن رموز لغة الحساب؛ وإن اخترناه ليذكّر بالضرب،
فليس $\times$ عملية الضرب نفسها، بل رمز دالة ثنائية
المحل، واسمه الرسمي $\Obj f^2_1$. وأما $\cdot$
فهو \emph{دالة الضرب} المعتادة. فـ$n\cdot m$ حاصل
ضرب العددين $n$ و~$m$؛ وأما $x\times y$ فحد من
حدود لغة الحساب. وفي النموذج القياسي نفسر~$\times$
بالدالة $\cdot$ على الأعداد الطبيعية. وأما $+$
فنستعمله للرمز اللغوي ولعملية الجمع نفسها؛ والسياق
هو الذي يميز المراد في كل موضع.

\begin{defn}
نسمي \emph{الحساب الصادق} نظرية ال!!{sentence}s التي يحققها
النموذج القياسي للحساب؛ فهي
\[
\Th{TA} = \Setabs{!A}{\Sat{N}{!A}}.
\]
\end{defn}

وكون $\Th{TA}$ نظرية يظهر مما يأتي: إذا كان
$\Th{TA}\Entails !A$، صدقت $!A$ في كل !!{structure}
تحقق~$\Th{TA}$. ومن $\Sat{N}{\Th{TA}}$ يلزم
$\Sat{N}{!A}$، فتكون $!A\in\Th{TA}$.

وأما الطريق الثاني إلى تعيين نظرية~$\Gamma$ فهو جمع
ال!!{sentence}s اللازمة من مجموعة جمل~$\Gamma_0$.
فتكون $\Gamma$ «إغلاق» $\Gamma_0$ تحت الاستلزام.
وإنما يفيد هذا الطريق إذا عُيّنت $\Gamma_0$ تعيينًا
صريحًا، كأن تسرد !!{element}sها؛ وأقل ما نطلبه أن
تقبل $\Gamma_0$ القرار، أي أن نستطيع باختبار حسابي
تمييز !!a{sentence} المنتمية إلى~$\Gamma_0$ من غيرها.
وتسمى ال!!{sentence}s في~$\Gamma_0$ \emph{بديهيات}
لـ~$\Gamma$، وتوصف $\Gamma$ بأنها \emph{مؤكسمة}
بـ~$\Gamma_0$.

\begin{defn}
نقول إن النظرية~$\Gamma$ \emph{مؤكسمة} بـ~$\Gamma_0$ متى، وفقط متى،
\[
\Gamma = \Setabs{!A}{\Gamma_0 \Entails !A}.
\]
\end{defn}

\begin{defn}
النظرية $\Th{Q}$ التي بديهياتها الجمل الآتية تسمى
«$\Th{Q}$ لروبنسون»، وهي من أبسط نظريات الحساب:
\begin{align*}
& \lforall[x][\lforall[y][(\eq[x'][y'] \lif \eq[x][y])]] \tag{$!Q_1$}\\
& \lforall[x][\eq/[\Obj 0][x']] \tag{$!Q_2$}\\
& \lforall[x][(\eq[x][\Obj 0] \lor \lexists[y][\eq[x][y']])] \tag{$!Q_3$}\\
& \lforall[x][\eq[(x + \Obj 0)][x]] \tag{$!Q_4$}\\
& \lforall[x][\lforall[y][\eq[(x + y')][(x + y)']]] \tag{$!Q_5$}\\
& \lforall[x][\eq[(x \times \Obj 0)][\Obj 0]] \tag{$!Q_6$}\\
& \lforall[x][\lforall[y][\eq[(x \times y')][((x \times y) + x)]]] \tag{$!Q_7$}\\
& \lforall[x][\lforall[y][(x < y \liff \lexists[z][\eq[(z' + x)][y]])]] \tag{$!Q_8$}
\end{align*}
فال!!{sentence}s في $\{!Q_1,\dots,!Q_8\}$ بديهيات $\Th{Q}$؛
وتجمع $\Th{Q}$ ال!!{sentence}s اللازمة منها جميعًا:
\[
\Th{Q} = \Setabs{!A}{\{!Q_1, \dots, !Q_8\} \Entails !A}.
\]
\end{defn}

\begin{defn}
لتكن $!A(x)$ !!a{formula} من $\Lang L_A$، ذوات المتغيرات
الحرة~$x$ و$y_1$، و\dots، و$y_n$. فكل !!{sentence} صورتها
\[
\lforall[y_1][\dots\lforall[y_n][((!A(\Obj 0) \land \lforall[x][(!A(x)
\lif !A(x'))]) \lif \lforall[x][!A(x)])]]
\]
نسخة من \emph{مخطط الاستقراء}.

ونسمي \emph{حساب بيانو}~$\Th{PA}$ النظرية التي تؤكسَم
ببديهيات $\Th{Q}$، مضافًا إليها جميع نسخ مخطط الاستقراء.
\end{defn}

\begin{explain}
كل نسخة من مخطط الاستقراء تصدق في~$\Struct{N}$.
وليتضح السبب، نفرض أن ال!!{formula}~$!A$ لا
!!{variable} حرًا فيها سوى~$x$. فتعيّن $!A(x)$ مجموعة
جزئية~$X_{!A}$ من~$\Nat$ في~$\Struct{N}$؛
و$X_{!A}$ تجمع كل~$n\in\Nat$ الذي يحقق
$\Sat{N}{!A(x)}[s]$ حين $s(x)=n$. والنسخة المقابلة
من مخطط الاستقراء هي
\[
((!A(\Obj 0) \land \lforall[x][(!A(x) \lif !A(x'))]) \lif
  \lforall[x][!A(x)]).
\]
فإن صدق المقدم في~$\Struct{N}$، كان $0\in X_{!A}$،
وكان $n+1\in X_{!A}$ كلما كان $n\in X_{!A}$.
فمن $0\in X_{!A}$ ينتج $1\in X_{!A}$، ومن
$1\in X_{!A}$ ينتج $2\in X_{!A}$، وهكذا.
فـ$n\in X_{!A}$ لكل $n\in\Nat$؛ وهذا هو تحقق
$\lforall[x][!A(x)]$ في~$\Struct{N}$.
\end{explain}

إذن $\Th{Q}$ و$\Th{PA}$ نظريتان مؤكسَمتان.
وإنما السؤال عن قوتهما: أتثبت $\Th{PA}$ جميع
الحقائق عن~$\Nat$ التي تعبّر عنها~$\Lang L_A$؟
أي أتحسم بديهيات $\Th{PA}$ كل سؤال يمكن صوغه
في~$\Lang L_A$؟

وهذا بعينه سؤال: هل $\Th{PA}=\Th{TA}$؟
فـ$\Th{TA}$ تثبت، أي تحتوي، الحقائق كلها عن~$\Nat$،
وتحسم كل سؤال يصاغ في~$\Lang L_A$. وذلك أن $!A$
إذا كانت !!a{sentence} في $\Lang L_A$، فإما
$\Sat{N}{!A}$ وإما $\Sat{N}{\lnot !A}$؛ ومن ثم
إما $\Th{TA}\Entails !A$ وإما $\Th{TA}\Entails\lnot !A$.
ونسمي النظرية التي هذا شأنها \emph{!!{complete}}.

\begin{defn}
النظرية $\Gamma$ \emph{!!{complete}} متى، وفقط متى، تحقق لكل
!!{sentence}~$!A$ في لغتها، إما $\Gamma\Entails !A$ وإما
$\Gamma\Entails\lnot !A$.
\end{defn}

\begin{explain}
وبمبرهنة التمام، $\Gamma\Entails !A$ تكافئ
$\Gamma\Proves !A$. فتكون $\Gamma$ تامة إذا وفقط إذا
كان لكل !!{sentence}~$!A$ من لغتها أحد الأمرين:
$\Gamma\Proves !A$ أو $\Gamma\Proves\lnot !A$.
\end{explain}

ونسأل أيضًا: أيصح اختبار انتماء !!a{sentence}
إلى~$\Th{TA}$ أو $\Th{PA}$، بل إلى~$\Th{Q}$ وحدها،
بإجراء حسابي؟ ولضبط هذا السؤال نحدّ متى تكون مجموعة،
ومنها مجموعات ال!!{sentence}s، !!{decidable}.

\begin{defn}
المجموعة $X$ \emph{!!{decidable}} إذا وفقط إذا أمكن
بإجراء حسابي، عند الدخل~$x$، إخراج~$1$ متى كان
$x\in X$، وإخراج~$0$ فيما عدا ذلك.
\end{defn}

فالسؤال إذن عن $\Th{TA}$، وكذلك $\Th{PA}$ و$\Th{Q}$:
أهي !!{decidable}؟

والجواب عن مسائل التقرير هذه كلها بالنفي؛ فلا واحدة
من هذه النظريات تقبل القرار. ولا يخصها ذلك وحدها،
بل يجري على \emph{كل} نظرية تستوفي شروطًا معينة.
ومن تلك الشروط، التي تحققها $\Th{Q}$ و$\Th{PA}$،
أن تكون مجموعة بديهياتها !!a{decidable}.

\begin{defn}
النظرية \emph{!!{axiomatizable}} هي ما أمكن أكسَمته بمجموعة
!!a{decidable} من البديهيات.
\end{defn}

\begin{ex}
إذا كانت بديهيات النظرية مجموعة منتهية من
ال!!{sentence}s، كانت !!{axiomatizable}؛ لأن المجموعة
المنتهية !!{decidable}. ومن أمثلة ذلك $\Th{Q}$،
فإنها !!{axiomatizable}.

وكذلك النظريات التي تعيّنها مخططات فعلية كهذه،
ومنها~$\Th{PA}$، !!{axiomatizable}. ولتمييز كون $!B$
من بديهيات~$\Th{PA}$، أي لحساب $\Char{X}$ حيث
$\Char{X}(!B)=1$ إن كانت $!B$ بديهية في~$\Th{PA}$،
و$=0$ خلاف ذلك، نسلك هذا السبيل: نختبر أولًا أ$!B$
إحدى بديهيات~$\Th{Q}$؛ فإن كانت، أجبنا بنعم وكانت
$\Char{X}(!B)=1$. وإلا فاختبار كونها نسخة من مخطط
الاستقراء ممكن على نظام. فإن كل نسخة تبدأ بكمّات كلية،
تليها !!{formula} فرعية شرطية، وتالي الشرطية
$\lforall[x][!A(x,y_1,\dots,y_n)]$. وفيها $x$ و$y_1$،
و\dots، و$y_n$ متغيرات~$!A$ الحرة كلها، والكمّات
الأولى في~$!B$ تربط~$y_1$، و\dots، و~$y_n$.
فنستخرج~$!A$، ونطابق متغيراتها الحرة مع المتغيرات
المربوطة في أول الجملة وبـ~$\lforall[x]$،
ثم نختبر أكان مقدم الشرطية هو:
\[
!A(\Obj 0, y_1, \dots, y_n) \land \lforall[x][(!A(x, y_1, \dots, y_n)
\lif !A(x', y_1, \dots, y_n))].
\]
فإن صحت المطابقة كانت $!B$ نسخة من المخطط؛ وإلا خرجت عنه.
\end{ex}

وفي بحث هذه المسألة، وما هو أعم منها من تمام النظريات
وقابليتها للقرار، ينفع التعريف الآتي. فقد تقدم أن $X$
!!{enumerable} إذا وفقط إذا كانت خالية أو وجدت
!!a{surjective} دالة~$f\colon\Nat\to X$؛ وهذه الدالة
تعداد لـ~$X$.

\begin{defn}
نقول عن المجموعة $X$ إنها \emph{!!{computably enumerable}} (واختصارًا
!!{c.e.}) متى كانت خالية أو أمكن تعدادها بدالة قابلة للحساب، وفقط حينئذ.
\end{defn}

ولا تكفي !!{axiomatizability} وحدها لانطباق مبرهنات
عدم الاكتمال؛ فمن الشروط أيضًا قدرة النظرية على إثبات
الحقائق الأساسية للدوال القابلة للحساب والعلاقات
!!{decidable}. والمراد بتلك الحقائق ال!!{sentence}s
التي تعيّن قيم الدوال عند كل مدخل، وبالقدرة أن تكون
هذه الجمل مبرهنات في النظرية. خذ الجمع مثالًا:
قيمة $+$ عند $2$ و$3$ هي~$5$، أي $2+3=5$.
والتعبير الحسابي عن أن قيمة $+$ عند $2$ و$3$
هي~$5$ هو $(\num{2}+\num{3})=\num{5}$؛ وهذه
الجملة تثبتها $\Th{Q}$. ونطلب على العموم لكل
دالة قابلة للحساب $f(x_1,x_2)$
!!a{formula}~$!A_f(x_1,x_2,y)$ في~$\Lang{L_A}$
تحقق $\Th{Q}\Proves !A_f(\num{n_1},\num{n_2},\num{m})$
كلما كان $f(n_1,n_2)=m$. فبذلك تثبت $\Th{Q}$
قيمة~$f$ عند $n_1$ و$n_2$، وأنها~$m$.
بل نزيد فنطلب إثبات انفراد هذه القيمة، فلا يكون
لـ~$f$ عند $n_1$ و~$n_2$ عدد آخر قيمةً.
ونظير ذلك في العلاقات !!{decidable}.
والتعريفان الآتيان يضبطان المراد.

\begin{defn}
نقول إن ال!!{formula}~$!A(x_1,\dots,x_k,y)$ \emph{!!{represents}} الدالة
$f\colon\Nat^k\to\Nat$ في~$\Gamma$ متى، وفقط متى، استلزم
$f(n_1,\dots,n_k)=m$ تحقق الشرطين:
\begin{enumerate}
\item $\Gamma \Proves !A(\num{n_1}, \dots, \num{n_k}, \num{m})$، و
\item $\Gamma \Proves \lforall[y](!A(\num{n_1}, \dots, \num{n_k},
y) \lif y = \num{m})$.
\end{enumerate}
\end{defn}

\begin{defn}
نقول إن ال!!{formula}~$!A(x_1,\dots,x_k)$ \emph{!!{represents}} العلاقة
$R\subseteq\Nat^k$ في~$\Gamma$ متى، وفقط متى، تحقق ما يأتي:
\begin{enumerate}
\item إذا تحقق $R(n_1,\dots,n_k)$، كان
$\Gamma\Proves !A(\num{n_1},\dots,\num{n_k})$، و
\item إذا لم يتحقق $R(n_1,\dots,n_k)$، كان $\Gamma\Proves\lnot
!A(\num{n_1},\dots,\num{n_k})$.
\end{enumerate}
\end{defn}

فالقدر من القوة المقصود لانطباق مبرهنات عدم الاكتمال
أن النظرية !!{represents} جميع الدوال القابلة للحساب
والعلاقات !!{decidable}. و$\Th{Q}$ وامتداداتها
تفي بهذا الشرط؛ إلا أن إثباته يحتاج إلى عمل،
لأنه حكم غير يسير على ما يمكن لـ$\Th{Q}$ إثباته.
ووجه الصعوبة أن بديهيات $\Th{Q}$ قليلة،
وعلينا استخراج الحقائق المطلوبة منها جميعًا.
لكن $\Th{Q}$، مع ذلك، نظرية ضعيفة جدًا.
فإثبات تمثيل $\Th{Q}$ للدوال القابلة للحساب
صعب، وأكثر النظريات المهمة أقوى من~$\Th{Q}$،
أي تزيد عليها فيما تثبت. وكل ما تثبته
$\Th{Q}$ تثبته النظرية الأقوى؛ فمتى ثبت
أن $\Th{Q}$ تمثل تلك الدوال، ثبت الأمر نفسه
لكل ما فوقها. لذلك ينتفع بهذا الجهد في طائفة
واسعة من النظريات، لا في نظرية واحدة.
بل أي نظرية تروم صوغ «الرياضيات كلها» ينبغي
أن تثبت على الأقل ما تثبته $\Th{Q}$،
إذ لا غنى لها عن استيعاب نتائج الحسابات
الأولية. فهي من هذه الجهة تدخل في نطاق
النظريات التي تتناولها مبرهنات عدم الاكتمال.

\end{document}
